% =============================================================================
%  Resumen Capítulo 3 — REGRESIÓN LINEAL
%  Libro: Evsukoff, A. (2020). Inteligência Computacional: Fundamentos e
%         Aplicações. Río de Janeiro: E-papers.
% =============================================================================
% =============================================================================
%  preambulo.tex — Preámbulo de los resúmenes en español
%  Evsukoff, A. (2020). Inteligência Computacional. Rio de Janeiro: E-papers.
% =============================================================================
\documentclass[11pt,a4paper]{article}

\usepackage[T1]{fontenc}
\usepackage[utf8]{inputenc}
\usepackage[spanish,es-tabla]{babel}
\usepackage{lmodern}
\usepackage[a4paper,margin=2.4cm]{geometry}
\usepackage{amsmath,amssymb,amsthm,mathtools}
\usepackage{bm}
\usepackage{enumitem}
\usepackage{xcolor}
\usepackage{tcolorbox}
\tcbuselibrary{skins,breakable}
\usepackage{hyperref}
\usepackage{microtype}
\usepackage{titlesec}
\usepackage{fancyhdr}
\usepackage{booktabs}
\usepackage{array}

\hypersetup{
  colorlinks=true,
  linkcolor=black!70,
  urlcolor=blue!60!black,
  citecolor=black!70,
  pdfborder={0 0 0}
}

\titleformat{\section}{\Large\bfseries\sffamily\color{black!85}}{\thesection}{0.7em}{}
\titleformat{\subsection}{\large\bfseries\sffamily\color{black!75}}{\thesubsection}{0.6em}{}

\pagestyle{fancy}
\fancyhf{}
\fancyhead[L]{\small\itshape Resumen — Inteligencia Computacional (Evsukoff, 2020)}
\fancyhead[R]{\small\thepage}
\renewcommand{\headrulewidth}{0.3pt}

\newtcolorbox{conceito}[1][]{
  colback=blue!4, colframe=blue!55!black, boxrule=0.6pt,
  arc=2mm, left=6pt, right=6pt, top=4pt, bottom=4pt,
  title={\bfseries Concepto clave}, fonttitle=\sffamily\bfseries,
  breakable, #1
}
\newtcolorbox{formula}[1][]{
  colback=gray!6, colframe=gray!50, boxrule=0.4pt,
  arc=1.5mm, left=6pt, right=6pt, top=3pt, bottom=3pt,
  breakable, #1
}
\newtcolorbox{aplicacao}[1][]{
  colback=green!4, colframe=green!45!black, boxrule=0.5pt,
  arc=2mm, left=6pt, right=6pt, top=3pt, bottom=3pt,
  title={\bfseries Aplicación}, fonttitle=\sffamily\bfseries,
  breakable, #1
}

\newcommand{\R}{\mathbb{R}}
\newcommand{\E}{\mathbb{E}}
\newcommand{\Var}{\operatorname{Var}}
\newcommand{\Cov}{\operatorname{Cov}}
\newcommand{\diag}{\operatorname{diag}}
\newcommand{\argmin}{\operatorname*{arg\,min}}
\newcommand{\argmax}{\operatorname*{arg\,max}}
\newcommand{\T}{^{\!\top}}
\DeclareMathOperator{\sign}{sign}
\DeclareMathOperator{\tr}{tr}


\title{\sffamily\bfseries Capítulo 3 — Regresión Lineal\\
       \large Mínimos cuadrados, regularización y validación de modelos}
\author{Resumen de estudio}
\date{}

\begin{document}
\maketitle
\thispagestyle{fancy}

\begin{conceito}
La \textbf{regresión lineal} ajusta un modelo
$\hat{y}(t)=f(\mathbf{x}(t);\boldsymbol\theta)$ que relaciona variables
de entrada $\mathbf{x}\in\R^p$ con una variable de salida numérica
$y\in\R$. ``Lineal'' se refiere a la \emph{linealidad en los
parámetros}, no necesariamente en los regresores: el modelo polinomial y
las funciones de base también caen en esta categoría. Es el primer
modelo a ajustar en cualquier aplicación --- simple, barato,
interpretable y, en muchos problemas reales, su diferencia frente a un
modelo no lineal es inferior al 5\%.
\end{conceito}

\section{El problema de regresión y el modelo lineal}

La relación real $y=f(\mathbf{x})$ es desconocida; se observan pares
$(\mathbf{x}(t),y(t))$ contaminados por ruido. El modelo aproxima esa
relación por
\begin{formula}
\[
\hat{y}(t)=\hat{\mathbf{x}}(t)\,\boldsymbol\theta\T
=\sum_{i=1}^{n}\hat{x}_i(t)\,\theta_i,
\]
donde $\hat{\mathbf{x}}(t)\in\R^n$ es el vector de \textbf{regresores} y
$\boldsymbol\theta=(\theta_1,\dots,\theta_n)$ los parámetros.
\end{formula}

Casos típicos:
\begin{itemize}[leftmargin=1.2cm]
  \item \textbf{Lineal simple (1\textsuperscript{er} orden):}
        $\hat{\mathbf{x}}=[1,\mathbf{x}]$, el modelo es un hiperplano,
        $n=p+1$.
  \item \textbf{Polinomial de grado $r$:}
        $\hat{\mathbf{x}}=[1,\mathbf{x},\mathbf{x}^2,\dots,\mathbf{x}^r]$,
        $n=r p+1$ despreciando términos cruzados.
  \item \textbf{Funciones de base $h_i(\mathbf{x})$:} permiten
        cualquier no linealidad fija en las entradas, manteniendo el
        ajuste lineal en los $\theta_i$ (Capítulo~7).
\end{itemize}

El residuo en el registro $t$ es $e(t)=y(t)-\hat{y}(t)$. El
\emph{overfitting} (sobreajuste) ocurre cuando la complejidad del
modelo excede lo que la muestra soporta: el modelo pasa a representar
el ruido. La Figura~3.4 del libro muestra cómo un polinomio de grado~10
ajustado a $N=20$ puntos generados por $y=x^2+\varepsilon$ adquiere
parámetros de magnitud absurda (hasta $\sim 130$) y pierde capacidad
de generalización.

\section{Método de Mínimos Cuadrados (MMC)}

La función de evaluación clásica es el \textbf{error cuadrático medio}
\[
\mathrm{ECM}(\boldsymbol\theta) =
\frac{1}{N}\sum_{t=1}^{N}\bigl(y(t)-\hat{y}(t)\bigr)^2
= \frac{1}{N}\bigl(\mathbf{Y}-\hat{\mathbf{X}}\boldsymbol\theta\T\bigr)\T
        \bigl(\mathbf{Y}-\hat{\mathbf{X}}\boldsymbol\theta\T\bigr).
\]
Bajo la hipótesis de residuos con distribución normal, media nula y
varianza constante (estacionariedad), \textbf{minimizar el ECM
equivale a maximizar la verosimilitud}; el MMC entrega parámetros de
mínima varianza (Gauss--Markov).

Anulando $\partial \mathrm{ECM}/\partial\boldsymbol\theta=
-\tfrac{2}{N}\hat{\mathbf{X}}\T(\mathbf{Y}-\hat{\mathbf{X}}\boldsymbol\theta\T)$
se llega a las \textbf{ecuaciones normales}
\begin{formula}
\[
\hat{\mathbf{X}}\T\hat{\mathbf{X}}\,\boldsymbol\theta\T
= \hat{\mathbf{X}}\T\mathbf{Y},
\quad\Longrightarrow\quad
\boldsymbol\theta^{*}=
\bigl(\hat{\mathbf{X}}\T\hat{\mathbf{X}}\bigr)^{-1}\hat{\mathbf{X}}\T\mathbf{Y},
\]
con \emph{matriz pseudoinversa}
$\hat{\mathbf{X}}^{+}=(\hat{\mathbf{X}}\T\hat{\mathbf{X}})^{-1}
\hat{\mathbf{X}}\T$ y \emph{matriz sombrero}
$\mathbf{H}=\hat{\mathbf{X}}\hat{\mathbf{X}}^{+}$ tal que
$\hat{\mathbf{Y}}=\mathbf{H}\mathbf{Y}$.
\end{formula}

\subsection{Interpretación geométrica}

Cada columna de $\hat{\mathbf{X}}$ es un vector en $\R^N$; las $n$
columnas generan un \emph{subespacio del modelo}. El MMC \textbf{proyecta
ortogonalmente} $\mathbf{Y}$ sobre ese subespacio, dejando el residuo
ortogonal al modelo:
\[
\hat{\mathbf{Y}}\T(\mathbf{Y}-\hat{\mathbf{X}}\boldsymbol\theta^{*})=\mathbf{0}.
\]
Consecuencia: la dimensión del modelo debe cumplir $n<N$, en caso
contrario el modelo es más complejo que la muestra. La \textbf{maldición
de la dimensionalidad} (Bellman) hace que, en alta dimensión, el número
de muestras necesarias para una cobertura razonable crezca
exponencialmente.

\section{Sesgo, varianza y regularización}

\subsection{Descomposición sesgo--varianza}

Para una entrada fija $\mathbf{x}$ con valor real $y_r=\E(y)$,
\[
\E\bigl((y_r-\hat{y})^2\bigr)
= \underbrace{\bigl(y_r-\E(\hat{y})\bigr)^2}_{\text{sesgo}^2}
+ \underbrace{\E\bigl((\hat{y}-\E(\hat{y}))^2\bigr)}_{\text{varianza}}.
\]
\begin{itemize}[leftmargin=1.2cm]
  \item Modelo \textbf{demasiado simple}: \emph{sesgo alto}, varianza
        baja (\emph{underfitting}).
  \item Modelo \textbf{demasiado complejo para la muestra}: sesgo bajo,
        \emph{varianza alta} (\emph{overfitting}).
\end{itemize}
La regularización busca \textbf{cambiar un poco de sesgo por mucha
reducción de varianza}, mejorando el error esperado.

\subsection{Regularización de Tikhonov (\emph{ridge regression})}

Se penaliza la norma del vector de parámetros minimizando
\begin{formula}
\[
\mathrm{ECMR}(\mu,\boldsymbol\theta)=
\frac{1}{N}\bigl\|\mathbf{Y}-\hat{\mathbf{X}}\boldsymbol\theta\T\bigr\|^2
+\mu\,\|\boldsymbol\theta\|^2,
\quad\mu\geq 0.
\]
La solución es
$\boldsymbol\theta_{\mu}^{*}=
\bigl(\hat{\mathbf{X}}\T\hat{\mathbf{X}}+\mu\mathbf{I}\bigr)^{-1}
\hat{\mathbf{X}}\T\mathbf{Y}$.
\end{formula}

El efecto es \textbf{numérico}: el número de condición de la matriz de
los coeficientes cambia de $\kappa=\lambda_{\max}/\lambda_{\min}$ a
$(\lambda_{\max}+\mu)/(\lambda_{\min}+\mu)$, garantizando estabilidad
incluso cuando $\hat{\mathbf{X}}\T\hat{\mathbf{X}}$ es singular o casi
singular (regresores muy correlacionados). Hoerl, Kennard y Baldwin
sugieren
$\mu=\tfrac{n}{N-n}\,\|\mathbf{Y}-\hat{\mathbf{X}}\boldsymbol\theta\T\|^2/
\|\boldsymbol\theta\|^2$.

\subsection{Regularización por componentes principales}

Si $\hat{\mathbf{X}}=\mathbf{U}\mathbf{S}\mathbf{V}\T$ es la SVD,
\[
\boldsymbol\theta_{r}^{*}=
\sum_{i=1}^{r}\frac{\mathbf{u}_i\T\mathbf{Y}}{s_i}\mathbf{v}_i,
\quad r\leq \mathrm{rank}(\hat{\mathbf{X}}).
\]
Truncar en los $r$ mayores valores singulares \textbf{descarta las
direcciones inestables} (asociadas a $s_i$ minúsculos), produciendo
una solución de menor norma y mejor condicionada. Equivale a aplicar el
MMC sólo en las componentes principales de la matriz de regresores.

\section{Validación de modelos}

\subsection{Estadísticas sobre el residuo}

La descomposición $\mathrm{SST}=\mathrm{SSR}+\mathrm{SSRes}$ separa
variabilidad total, explicada y residual. Se definen:
\begin{align*}
R^2 &= 1-\frac{\mathrm{SSRes}}{\mathrm{SST}}
     = 1-\frac{\sum (y(t)-\hat{y}(t))^2}{\sum(y(t)-\bar{y})^2}, \\
\mathrm{RMS} &= \sqrt{\tfrac{1}{N}\textstyle\sum e^2(t)}, \\
\mathrm{MAE} &= \tfrac{1}{N}\textstyle\sum |e(t)|, \\
\mathrm{MAPE} &= \tfrac{1}{N}\textstyle\sum |e(t)/y(t)|, \\
\rho_{y\hat{y}} &= \mathrm{Corr}(y,\hat{y}).
\end{align*}
$R^2\approx 1$ es deseable; $R^2\leq 0$ indica que el modelo es peor
que la media. Para variables con gran amplitud, usar $R^2_{\log}$.

La \textbf{varianza de los residuos} se estima por
$\hat\sigma^2=\tfrac{1}{N-n}\sum e^2(t)$ (denominador $N-n$ por grados
de libertad). Bajo normalidad, intervalo de confianza al 95\% para la
predicción en $\hat{\mathbf{x}}_0$:
\[
\hat{y}_0 \pm 2\sqrt{\hat\sigma^2 \hat{\mathbf{x}}_0\T
(\hat{\mathbf{X}}\T\hat{\mathbf{X}})^{-1}\hat{\mathbf{x}}_0}.
\]
La normalidad de los residuos se diagnostica por histograma y gráfico
de probabilidad normal.

\subsection{Validación cruzada}

Calcular las estadísticas en el propio conjunto de entrenamiento
favorece modelos con \emph{overfitting}. Criterios penalizados como
AIC y BIC ayudan pero no siempre aciertan. La solución práctica es la
\textbf{validación cruzada de $K$ ciclos} (Algoritmo~3.1):

\begin{enumerate}[leftmargin=1.2cm]
  \item Particionar $\mathcal{T}$ en $K$ subconjuntos
        aproximadamente iguales.
  \item Para cada $i=1,\dots,K$: ajustar el modelo en
        $\mathcal{T}\setminus\mathcal{T}_i$ y evaluar en
        $\mathcal{T}_i$.
  \item Concatenar las predicciones $\hat{\mathbf{Y}}_i$ y calcular las
        estadísticas globales.
\end{enumerate}

Casos especiales: \emph{leave-one-out} ($K=N$) es el estimador más
preciso pero el más caro y sensible a la muestra; valores típicos
$K=5$ o $K=10$ ofrecen buen compromiso.

\section{Aplicaciones}

\begin{aplicacao}
\textbf{Pollution} (contaminación del aire y mortalidad, EE.UU.,
1959--1961, 60~registros, 15~variables). La eliminación de
\emph{outliers} (12\%) eleva $R^2$ de 0{,}41 a 0{,}60. La regularización
de Tikhonov con $\mu=0{,}01$ sobre los datos sin \emph{outliers} produce
el mejor modelo: $R^2=0{,}61$, $\mathrm{RMS}=0{,}206$,
$\mathrm{MAPE}=18{,}5\%$. La combinación ``eliminación de
\emph{outliers} + regularización'' confirma la intuición de la
descomposición sesgo--varianza.
\end{aplicacao}

\begin{aplicacao}
\textbf{Elevación adiabática de temperatura del hormigón}
(Eletrobrás Furnas, 136~registros, 14~entradas físico-químicas, 3
salidas paramétricas). La variable $y_1$ (temperatura final) es la más
fácil ($R^2=0{,}78$); $y_3$ es la más difícil ($R^2=0{,}39$). La
regularización mejora ligeramente todos los indicadores. El problema
es intrínsecamente no lineal; modelos basados en redes neuronales y
sistemas \emph{fuzzy} (Capítulos~8--9) entregan predicciones más
precisas.
\end{aplicacao}

\section{Comentarios finales}

\begin{enumerate}[leftmargin=1.2cm]
  \item Modelar es \textbf{elegir}: estructura $\rightarrow$ ajuste de
        parámetros $\rightarrow$ validación. La elección de la
        estructura es la etapa de mayor impacto.
  \item El MMC es elegante y óptimo bajo sus hipótesis, pero
        inestable cuando $n$ es grande comparado con $N$; la
        \textbf{regularización} (Tikhonov o componentes principales)
        es la defensa estándar.
  \item Siempre validar en datos \textbf{no vistos} (validación
        cruzada). Las estadísticas sobre el conjunto de entrenamiento
        pueden enmascarar \emph{overfitting}.
  \item La eliminación de \emph{outliers} puede mejorar mucho el
        ajuste, pero debe aplicarse con cautela: se descartan
        justamente los registros más difíciles.
  \item El capítulo siguiente (4) extiende este arsenal a las
        \emph{series temporales} y los \emph{sistemas dinámicos
        lineales} (AR, MA, ARMA, ARX, ARMAX, OE, Box \& Jenkins,
        ecuaciones de estado).
\end{enumerate}

\end{document}
